\section{Proceso Formal de Gestión de Defectos}

El ciclo de vida de los defectos funcionales, lógicos y de seguridad se gestionará bajo un flujo formalizado que garantice que ningún \textit{bug} crítico afecte la experiencia en producción.

\subsection{Herramienta de Seguimiento}
Se adoptará \textbf{Linear App} como la herramienta única de registro, asignación y auditoría de \textit{bugs}.

\subsection{Campos Requeridos para el Registro de un Defecto}
Cada vez que el equipo de Aseguramiento de Calidad (o el Gestor Operativo) identifique un comportamiento anómalo, abrirá una ``Issue'' en Linear App completando de manera mandatoria la siguiente plantilla estructural:

\begin{itemize}
    \item \textbf{Título del Error:} Breve y descriptivo, iniciando con la etiqueta de severidad (ej. \texttt{[BUG - ALTA] El historial de estados no muestra la autoría del cambio}). En Linear se pueden usar etiquetas de prioridad para indicar la prioridad del error (Alta, Media, Baja).
    \item \textbf{Descripción del Comportamiento Actual:} Redacción textual del fallo y la inconsistencia detectada.
    \item \textbf{Comportamiento Esperado:} Detalle del comportamiento correcto según las Historias de Usuario o criterios de aceptación del plan.
    \item \textbf{Pasos Exactos para su Reproducción:} Lista numerada con las entradas, URLs, datos ingresados y clics realizados para recrear el error.
    \item \textbf{Evidencias Adjuntas:} Capturas de pantalla, archivos de log del servidor Laravel (\texttt{storage/logs/laravel.log}) o capturas de la consola de desarrollador del navegador.
    \item \textbf{Prioridad / Criticidad:} Clasificación formal en Crítica, Alta, Media o Baja.
    \item \textbf{Responsable de la Corrección:} Asignación explícita de un desarrollador del equipo para resolverlo.
    \item \textbf{Comunicación con el equipo:} El responsable de abrir la ``Issue'' deberá comunicar por los grupos de WhatsApp o Teams para notificar de la situación.
\end{itemize}

\subsection{Acuerdos de Nivel de Servicio (SLA) de Calidad}
De acuerdo a la criticidad del defecto, el equipo acuerda los siguientes tiempos límite de respuesta (primera atención/asignación) y resolución (defecto corregido y verificado en \textit{staging}):

\begin{table}[htbp]
\centering
\small
\begin{tabularx}{\textwidth}{l X c c}
\toprule
\textbf{Nivel de Criticidad} & \textbf{Criterio Técnico de Clasificación} & \textbf{SLA Resp.} & \textbf{SLA Resol.} \\
\midrule
\textbf{Crítica (Critical)} & 
Defecto que inhabilita completamente la operación del sistema sin \textit{workaround} posible (ej: caída de PostgreSQL, fallos en la autenticación general, errores 500 al guardar cualquier reporte). & 
$< 2$ horas & 
$< 8$ horas \\
\midrule
\textbf{Alta (High)} & 
Defecto que degrada significativamente un módulo crítico, impidiendo su objetivo de negocio, pero el sistema sigue en marcha (ej: mapas Leaflet que no cargan marcadores, o transiciones de estado de incidencias que fallan). & 
$< 4$ horas & 
$< 24$ horas \\
\midrule
\textbf{Media (Medium)} & 
Errores funcionales parciales o discrepancias en criterios de aceptación con \textit{workaround} operativo (ej: comentarios editables después de 30 minutos, problemas de paginación en listados). & 
$< 12$ horas & 
$< 72$ horas \\
\midrule
\textbf{Baja (Low)} & 
Defectos cosméticos, de alineación visual en Bootstrap o textos y ortografía, que no alteran en absoluto la operatividad funcional de la plataforma. & 
$< 24$ horas & 
5 días hábiles \\
\bottomrule
\end{tabularx}
\caption{Tiempos de respuesta y resolución según criticidad (SLA)}
\label{tab:sla_defectos}
\end{table}

\subsection{Ejemplificación de Registro de Defecto (Linear App)}
Para ilustrar la aplicación del estándar de reporte definido en la sección 12.2, se presenta a continuación el registro formal del defecto correspondiente a las inconsistencias detectadas en la pantalla de gestión:

\begin{itemize}
    \item \textbf{Título del Error:} \texttt{[BUG - ALTA] Inconsistencias visuales del listado y error de eliminación de incidencias recién creadas}
    \item \textbf{Descripción del Comportamiento Actual:}
    Al ingresar a la pantalla de Gestión de Incidencias (\texttt{http://localhost/html/vista.html}), se observan múltiples desviaciones estéticas y una falla funcional de criticidad alta:
    \begin{enumerate}
        \item \textit{Layout de filtros:} El panel superior de filtros se distribuye en dos filas, desperdiciando espacio vertical en la interfaz.
        \item \textit{Ausencia de búsqueda:} El selector ``Todas las incidencias'' carece de un campo de búsqueda incremental/predictiva.
        \item \textit{Desbordamiento de texto:} La columna ``Descripción'' muestra el texto completo de los reportes en lugar de limitarse a una sola línea, desordenando la altura de las filas.
        \item \textit{Formato de fuentes:} El contenido de las columnas ``Reportado por'' y ``Ubicación'' se renderiza en negrita (fuente de peso \textit{bold}).
        \item \textit{Disposición de columnas:} El orden visual de las columnas de la tabla no es el requerido.
        \item \textit{Falta de responsividad:} En resoluciones de pantalla estándar de $1366\times768$ o superiores, la visualización de las 8 columnas genera un scroll horizontal innecesario.
        \item \textit{Paginador estático:} Los botones del paginador y el indicador del total de registros tienen estilos Bootstrap inconsistentes y el selector numérico ``Mostrar X elementos'' no actualiza el listado.
        \item \textit{Falla en eliminación:} Si un usuario crea una incidencia y, sin refrescar el navegador, intenta eliminarla desde el listado, el registro permanece visible en la pantalla. Al hacer clic por segunda vez sobre el botón ``Eliminar'', se genera un error HTTP genérico sin causa descrita. El registro solo desaparece del DOM si se fuerza una acción de edición y posterior cancelación.
    \end{enumerate}
    
    \item \textbf{Comportamiento Esperado:}
    La interfaz debe alinearse con los criterios de usabilidad del plan y el backend debe responder de manera semántica:
    \begin{enumerate}
        \item \textit{Layout de filtros:} Los filtros deben compactarse y alinearse en una única fila horizontal.
        \item \textit{Búsqueda predictiva:} Incluir un cuadro de texto incremental que filtre dinámicamente por título o ID de la incidencia mediante llamadas asíncronas de Fetch API.
        \item \textit{Truncamiento de texto:} La columna ``Descripción'' debe limitar su texto a una sola línea (utilizando truncamiento CSS con elipsis \texttt{...}), desplegando el contenido total únicamente al posicionar el cursor sobre la celda mediante un tooltip estándar de Bootstrap.
        \item \textit{Formato de fuentes:} Las columnas ``Reportado por'' y ``Ubicación'' deben renderizarse con peso de fuente normal (\texttt{font-weight: normal}).
        \item \textit{Orden de columnas:} El listado debe estructurarse estrictamente en la secuencia: \textit{Código, Título, Reportado por, Tipo, Prioridad, Descripción, Caracteres, Estado}.
        \item \textit{Ajuste responsivo:} La tabla con las 8 columnas debe autoajustarse dinámicamente y ser completamente legible sin scroll horizontal en resoluciones iguales o superiores a $1366\times768$.
        \item \textit{Paginador dinámico:} El paginador final debe calcular de manera dinámica las páginas en base al total de incidencias devueltas por la API, mostrar estilos estándar de Bootstrap y reaccionar al cambio del selector de cantidad de registros.
        \item \textit{Flujo de eliminación seguro:} Al accionar el botón de eliminación sobre cualquier registro, se debe ejecutar una petición AJAX \texttt{DELETE} exitosa, remover inmediatamente la fila del DOM y, en caso de fallo relacional o de lógica en el servidor (ej: incidencia cerrada o con comentarios críticos vinculados), retornar un código de error HTTP estructurado y un mensaje informativo preciso en pantalla.
    \end{enumerate}
    
    \item \textbf{Pasos Exactos para su Reproducción:}
    \begin{enumerate}
        \item Iniciar sesión en la aplicación con perfil de Gestor Operativo.
        \item Acceder al listado general en \texttt{http://localhost/html/vista.html}.
        \item Validar visualmente el desbordamiento de la columna ``Descripción'', el scroll horizontal a $1366\times768$ y la doble fila de filtros.
        \item Ir al formulario e ingresar una nueva incidencia (``Incidencia de prueba de SQA'').
        \item Retornar al listado y hacer clic inmediatamente en el botón ``Eliminar'' de la incidencia agregada.
        \item Observar que la fila sigue visible en la tabla.
        \item Hacer clic en ``Eliminar'' nuevamente y revisar el error 500 arrojado en la pestaña de red de la consola del navegador.
    \end{enumerate}
    
    \item \textbf{Evidencias Adjuntas:}
    \begin{itemize}
        \item \textit{Log del Servidor Laravel:} Exception \texttt{QueryException} en PostgreSQL debido a violación de clave foránea en la tabla relacional al intentar ejecutar \texttt{DELETE} sin control de dependencias.
        \item \textit{Captura de Consola:} \texttt{DELETE http://localhost/api/incidencias/99 500 (Internal Server Error)}.
        \item \textit{Capturas de Pantalla UI:} Listado con scroll horizontal activo y distorsión de la altura de la fila por descripción larga.
    \end{itemize}
    \item \textbf{Prioridad / Criticidad:} Alta (degradación funcional en el CRUD de incidencias y discrepancias múltiples de usabilidad).
    \item \textbf{Responsable de la Corrección:} Desarrollador Frontend (\texttt{@dev-front}) + Desarrollador Backend (\texttt{@dev-back}).
    \item \textbf{Módulo Afectado:} Frontend (Gestión de Vistas) / Backend API (Endpoints de Incidencia).
    \item \textbf{Sprint Asociado:} Sprint 3 --- Correcciones UI/UX y estabilidad.
    \item \textbf{Estimación de Esfuerzo:} 6 horas de desarrollo y pruebas.
    \item \textbf{Comunicación con el equipo:} Error registrado en Linear App y notificado al desarrollador backend a través del canal técnico de Teams para priorizar la verificación de integridad referencial.
\end{itemize}